\textbf{Nutrition5k has no volume ground truth.} The dataset provides
overhead RealSense depth and weighed mass, not food volume. We therefore score
monocular models against the volume integrated from sensor depth over the same
dish region (food plus plate rim above the support plane). This measures how
well a model reproduces a depth sensor, not how well it measures food. Our
attempt to separate food from plate with a second plate-floor plane fit gave a
food volume that barely tracks weighed mass (Spearman $\rho=0.24$,
Pearson $r=-0.01$, $n=502$), so we do not use it as a target. Absolute food-volume
accuracy is reported only on synthetic scenes with exact mesh volume.

\textbf{Synthetic-to-real gap.} Blender scenes have perfect camera intrinsics,
controlled lighting and procedural textures. Real OV2640 frames add rolling
shutter, auto-exposure drift, lens distortion and JPEG blocking on low-texture
food. The resolution and JPEG ablations bound part of this gap but not all of it. Textureless food already fails in simulation (Sec.~\ref{sec:synth}); an active texture projector or image-space segmentation would be needed for it.

\textbf{No physical measurements yet.} Load-cell linearity, repeatability and
drift, turntable angular error, scan duration and Wi-Fi transfer latency have
scripts and empty tables (Sec.~\ref{sec:hw}); none of the numbers in this paper
come from the physical rig. Timings are CPU timings on 8-vCPU cloud VMs, not
on the ESP32-S3.

\textbf{Volume is not mass.} Converting volume to grams needs a per-food
density. We do not estimate density; the load cell measures mass directly, which
is the reason it is part of the rig.
